\chapter{Introduction}
\label{chap:introduction}

Music festivals and large outdoor events have become a massive part of modern culture. Every summer, millions of people gather in open fields, parks, and temporary venues to enjoy live performances and food. From Untold in Romania to Tomorrowland in Belgium, these events keep growing in size and complexity. But behind all the excitement, there is a very practical problem that organizers have to deal with: how do you let tens of thousands of people buy food, drinks, and merchandise in a place where the internet barely works?

The traditional approach was straightforward: people brought cash, stood in line, and paid at the counter. It was slow, it was messy with always having change and there was always the risk of theft or losing your wallet in the crowd. Over the past decade, many festivals have moved towards cashless payment systems, typically based on RFID wristbands or prepaid cards. The idea is simple: you load money onto your wristband before or during the event, and then you just tap it at a vendor terminal to pay. Industry data from cashless technology providers suggests that RFID systems reduce transaction times by up to 80 percent compared to cash and that festival organizers see average revenue increases of around 22 percent after going cashless \cite{Tappit2024}. These figures are consistent with the academic evidence: a recent meta-analysis across 71 studies found a statistically significant, small positive effect of cashless payments on consumer spending ($g = 0.135$, 95\% CI [0.068; 0.203]), driven primarily by the reduced psychological friction of not physically handing over money \cite{Schomburgk2024}.

However, these systems come with a critical weakness: most of them depend on a reliable internet connection. And if there is one thing that festivals are known for, it is terrible connectivity. When tens of thousands of people gather in the same area, all using their phones simultaneously, the cellular network simply cannot handle the load. Shafiq et al.~\cite{Shafiq2013} conducted the first large-scale characterization of cellular network performance during crowded events using real-world traces from a tier-1 US network. Their findings are impressive: pre-connection failures increased by 100 to 5,000 times compared to routine days, driven by users exhausting the limited bandwidth of the signaling channel when too many devices attempt to acquire radio resources simultaneously. This degradation extended as far as 10 miles around the event venue as attendees arrived and departed. Even for users who managed to connect, voice call failures increased by 7 to 30 times, packet, and round-trip times increased by up to 7 times. The performance spikes were particularly severe during intermissions, when large numbers of attendees simultaneously attempted to use their devices. Notably, the study found that even with portable base stations and free Wi-Fi offloading already deployed, these degradation levels persisted, highlighting a fundamental capacity limitation rather than a simple configuration problem.


This is where things get problematic for payment systems. If the payment terminal at a food stall needs to contact a remote server to authorize a transaction, and the network is down or painfully slow, the transaction fails. The attendee cannot buy food. The vendor loses a sale. Multiply this by tens of thousands of transactions per day, and you can see why this is a serious issue. The most famous example of this going wrong happened at the Download Festival in the UK in 2015, the first major UK festival to go completely cashless. The RFID system experienced technical failures that left attendees unable to purchase food or drinks, with many taking to social media to describe the system as useless. Although organizers claimed only around one percent of attendees were affected, the backlash was severe enough that the festival abandoned the cashless-only approach the following year \cite{eFestivals2016}.

The core question this thesis tries to answer is: can we build a payment system that works reliably at festivals and that does not fall apart when the network has a bad moment? Our approach is to design a system that operates online by default, with connectivity coming from a Starlink uplink and a redundant network topology, so that vendor terminals can talk to the cloud backend in real time most of the time. But no network is perfectly reliable, especially one assembled in a field for a few days, so we also need a fallback for the moments when the network is down. Instead of building a separate offline server or a heavy synchronization layer, we let the NFC wristband itself act as the cache, holding the balance and a small counter that the vendor terminal can update locally until the network comes back.

\section{Problem Statement}
\label{sec:problem-statement}

Festivals put together a few problems that, on their own, payment systems already know how to solve, but putting them all together is when it becomes tricky. The first one is connectivity. Cellular networks fall over when too many people gather in the same place, and many festivals happen in fields and forests where there is no proper wired infrastructure to begin with. Our answer is to bring our own connectivity to the venue: a Starlink uplink with a redundant network topology, so that the vendor terminals can talk to the cloud backend most of the time. When the network is up, transactions are authorized in real time, wallet balances are fresh, and operations that need an external provider like Stripe work as you would expect.

A network we set up ourselves is still a network, and it can still go down. Equipment fails, the density of wireless devices in a small area creates dead zones, and multi-day events keep stressing infrastructure that was put together in a field a few days before. When the network does drop for a few minutes during a peak hour, the system still has to take payments, otherwise people cannot buy a sandwich and the whole festival economy stops. So we need a fallback that takes over when the network is gone and gets out of the way when it comes back. We do not want to design the whole system around the assumption that the network will fail, because that would push complexity everywhere for a problem that is online most of the time.

The second piece is that this is a closed economy. Unlike a normal shop where you tap a bank card and the bank takes care of authorization, a festival has its own internal currency: attendees load tokens, spend them at approved vendors, and get refunded for what they did not use. The money never leaves the perimeter of the event. This is actually an advantage for us, because we control the whole flow and we are not blocked by an external authorization network that we could not reach during an outage anyway.

The third piece is scale and speed. A medium festival can have hundreds of vendors processing thousands of transactions per hour, with very sharp peaks around meals and between concerts. Nobody is going to wait two minutes for a payment to go through. Online, this is just a matter of provisioning the backend to handle the spike. Offline, the question becomes how to keep the same speed at each terminal without giving up correctness, which is the next problem.

The fourth piece is consistency during offline windows. If two terminals lose the network at the same time and both approve a payment from the same wallet, they could together let through more money than is actually there. This is the well-known double-spend problem, and it is pretty hard in distributed systems. Refusing the payment is not an option for us, because at a festival "we cannot pay because the Wi-Fi is bad" is not an acceptable answer. So we choose availability, and we make sure the architecture limits how badly things can go wrong: each wristband carries its own balance and a debit counter on the chip itself, so the chip is the one that decides locally whether a payment fits, and double-spend across terminals is not possible because the chip is a single physical object that the customer carries between vendors. If minor discrepancies arise the festival will cover the loss, with the maximum exposure bounded by the queue size cap times the maximum transaction amount per terminal.

The fifth piece is reconciliation. When terminals come back online, the transactions they processed offline have to land in the server in a way that does not lose anything and does not double-apply anything. Because we are online-first, these windows are expected to be short and not that frequent, but even a few minutes of outage during a peak can produce hundreds of queued transactions per terminal, and the sync has to be idempotent and order-tolerant.

One last thing worth saying out loud: not every operation can work offline. Wallet top-ups in particular need real-time communication with Stripe (or any other on-ramp provider) and with our backend, because real money is moving from the attendee's bank account into the system. You cannot process that on the chip. This is on purpose: the chip handles debits when the network is gone, but credits, refunds, and anything that touches external money stay strictly online and stay strictly on the server.

\section{Objectives}
\label{sec:objectives}

The main objective of this thesis is to design and implement a minimum viable product (MVP) for a cloud-based, closed-loop payment system tailored to festival and event environments. The system is online by default, with connectivity coming from a Starlink uplink and a redundant network topology, and uses the NFC wristband itself as a local cache so that vendor terminals can keep approving payments during short network outages. The concept of a minimum viable product, as defined by Ries~\cite{Ries2011}, refers to a version of a product that allows the team to collect the maximum amount of validated learning with the least effort. In our case, this means building enough of the system to demonstrate that the core architectural ideas work, without necessarily implementing every feature a production deployment would need. Throughout this work, we explore different design architectures and patterns, discussing their benefits and trade-offs in the context of festival environments, before settling on the approach that best fits the constraints of the problem.

Specifically, the objectives of this work are the following:

\textbf{O1:} Design a payment architecture that keeps a closed-loop wallet usable through short network outages, without splitting the codebase into separate online and offline implementations. The same vendor-side flow has to work whether the network is up or down, with the chip and the server taking turns as the local source of truth depending on what is reachable.

\textbf{O2:} Build the wallet service that supports the two operations a festival needs in practice: top-ups through an external payment provider while online, since real money is entering the system, and debits at vendor terminals both online and during outages, with the terminal queueing offline debits for a later batch sync.

\textbf{O3:} Specify a reconciliation model that converges back to a consistent balance after offline windows without ad-hoc conflict resolution, and verify it with end-to-end tests that simulate the cases the model has to handle (idempotent retries, contention between two terminals, top-ups that race with offline debits).

\textbf{O4:} Implement a core service that manages vendors, products, stock, users, and events, providing the administrative backbone of the system.

\textbf{O5:} Build the client applications attendees, vendors, and event organizers need: a single mobile app that adapts to the attendee or vendor role on the same device, plus a separate web dashboard for organizers to set up and run events.

\textbf{O6:} Discuss the trade-offs of the chosen architecture and verify it against the festival requirements outlined in Section~\ref{sec:problem-statement}, deferred to Chapter~\ref{chap:discussion}.

\section{Thesis structure}
\label{sec:thesis-structure}

The rest of the thesis follows the path from the prior art to the running system and back to a verdict. Chapter~\ref{chap:background} reviews the existing work that informs the design: transit cards, commercial festival systems, EMV's offline authorisation, and the distributed systems literature on conflict resolution. Chapter~\ref{chap:architecture} presents the system architecture, with the hybrid two-counter reconciliation model at its core. Chapter~\ref{chap:implementation} walks through the parts of the running code that matter for that model: the database schema rework, the on-chip data layout, the sync endpoint, and the offline queue. Chapter~\ref{chap:security} discusses the threat model and the load-bearing weakness around cryptographic key management in the current MVP. Chapter~\ref{chap:discussion} returns to the trade-offs, compares the result against commercial systems, and evaluates the architecture against the objectives stated above. Chapter~\ref{chap:conclusions} summarises what the project did and did not achieve.
